\chapter*{Περίληψη}
\addcontentsline{toc}{chapter}{Περίληψη}

\guide{Γράφεται τελευταία, όταν υπάρχουν αριθμοί από το \cref{ch:evaluation}. 250--350 λέξεις, μία παράγραφος ή τέσσερις μικρές:
(1) πρόβλημα, (2) τι κατασκευάστηκε, (3) πώς αξιολογήθηκε, (4) τα κύρια αποτελέσματα \emph{με αριθμούς} (π.χ. «οριζόντιο σφάλμα \fillin{x}\,cm σε απόσταση \fillin{y}\,m από τα σημεία ελέγχου», «σταθεροί 72\,fps έως \fillin{z} χιλιάδες τρίγωνα», «SUS \fillin{w}»).}

Οι μετρήσεις των αισθητήρων παρακολούθησης κατασκευών απεικονίζονται συνήθως σε οθόνες, μακριά από το σημείο όπου λαμβάνονται, και ο μηχανικός στο πεδίο πρέπει να αντιστοιχίσει νοερά κάθε σημείο του χάρτη με τη θέση του στην κατασκευή. Η εργασία παρουσιάζει ένα end-to-end σύστημα που μεταφέρει τις μετρήσεις από μια πλατφόρμα παρακολούθησης IoT σε ένα αυτόνομο headset μικτής πραγματικότητας (Meta Quest~3) και τις εμφανίζει πάνω στους ίδιους τους αισθητήρες, στον πραγματικό χώρο.
\fillin{Συνέχεια: ευθυγράμμιση, απεικονίσεις, αξιολόγηση, αποτελέσματα.}

\vspace{1em}
\noindent\textbf{Λέξεις-κλειδιά:} μικτή πραγματικότητα, παρακολούθηση κατασκευών, διαδίκτυο των πραγμάτων, γεωαναφορά, χωρική ευθυγράμμιση, Meta Quest~3, οπτικοποίηση δεδομένων

\begin{otherlanguage}{english}
\chapter*{Abstract}
\addcontentsline{toc}{chapter}{Abstract}

\guide{The English abstract is a translation of the Greek one, not a different text.}

Measurements from structural monitoring sensors are usually viewed on screens, away from where they are taken, so an engineer on site has to map every point on the plan to its place on the structure. This thesis presents an end-to-end system that carries measurements from an IoT monitoring platform to a standalone mixed reality headset (Meta Quest~3) and shows them on the sensors themselves, in physical space.
\fillin{Continued.}

\vspace{1em}
\noindent\textbf{Keywords:} mixed reality, structural health monitoring, Internet of Things, georeferencing, spatial registration, Meta Quest~3, data visualization
\end{otherlanguage}
